\newcounter{exercice}
\newtcolorbox{exercicebox}[1][\unskip]{
    colback=BleuClair,
    colframe=Bleu,
    boxrule=1pt,
    arc=3mm,
    title=\bb{Exercice \stepcounter{exercice}\theexercice #1},
    left=5mm,
    right=5mm,
    top=2mm,
    bottom=2mm
}


\newcounter{probleme}
\newtcolorbox{probbox}[1][\unskip]{
    colback=RougeClair,
    colframe=Rouge,
    boxrule=1pt,
    arc=3mm,
    title=\bb{Problème \stepcounter{probleme}\theprobleme #1},
    left=5mm,
    right=5mm,
    top=2mm,
    bottom=2mm
}

\newenvironment{tableau}[1]
{%
    \renewcommand{\arraystretch}{1.3}%
    \begin{center}%
    \begin{tabularx}{\linewidth}{#1}%
    \toprule
}
{%
    \bottomrule
    \end{tabularx}%
    \end{center}%
    \renewcommand{\arraystretch}{1}%
}

